\documentclass[11pt,oneside]{article}

\usepackage[T1]{fontenc}
\usepackage{lmodern}
\usepackage{microtype}
\usepackage[margin=1in,headheight=14pt]{geometry}
\usepackage{amsmath,amssymb,amsthm,mathtools,mathrsfs}
\usepackage{booktabs,longtable,tabularx,array}
\usepackage{enumitem}
\usepackage{xcolor}
\usepackage{xurl}
\usepackage{hyperref}
\usepackage{fancyhdr}
\usepackage{titlesec}
\usepackage{listings}
\usepackage{needspace}

\definecolor{provedgreen}{RGB}{15,105,68}
\definecolor{conditionalblue}{RGB}{24,91,145}
\definecolor{opengray}{RGB}{85,92,99}
\definecolor{computeviolet}{RGB}{98,55,145}
\definecolor{codegray}{RGB}{248,248,248}
\definecolor{rulegray}{RGB}{92,101,111}

\hypersetup{
  colorlinks=true,
  linkcolor=conditionalblue,
  citecolor=provedgreen,
  urlcolor=conditionalblue,
  pdftitle={Algebraic Traps on the Power Curve: New Progress Toward the Alaoglu--Erdos Conjecture},
  pdfauthor={Research progress report},
  pdfsubject={Coppersmith--Chebyshev certificates, total positivity, and successive-minimum targets}
}
\urlstyle{same}

\pagestyle{fancy}
\fancyhf{}
\lhead{Alaoglu--Erd\H{o}s progress report V}
\rhead{Algebraic traps on the power curve}
\cfoot{\thepage}
\renewcommand{\headrulewidth}{0.4pt}

\titleformat{\section}{\Large\bfseries}{\thesection}{0.75em}{}
\titleformat{\subsection}{\large\bfseries}{\thesubsection}{0.75em}{}
\titleformat{\subsubsection}{\normalsize\bfseries}{\thesubsubsection}{0.75em}{}

\setlist[itemize]{leftmargin=2em,itemsep=0.28em,topsep=0.4em}
\setlist[enumerate]{leftmargin=2.3em,itemsep=0.28em,topsep=0.4em}
\setlength{\parindent}{0pt}
\setlength{\parskip}{0.55em}
\setlength{\emergencystretch}{3em}
\allowdisplaybreaks

\newtheorem{theorem}{Theorem}[section]
\newtheorem{proposition}[theorem]{Proposition}
\newtheorem{lemma}[theorem]{Lemma}
\newtheorem{corollary}[theorem]{Corollary}
\newtheorem{criterion}[theorem]{Criterion}
\newtheorem{target}[theorem]{Research target}
\newtheorem{conjecture}[theorem]{Conjecture}
\theoremstyle{definition}
\newtheorem{definition}[theorem]{Definition}
\newtheorem{experiment}[theorem]{Computational experiment}
\newtheorem{warning}[theorem]{Caution}
\theoremstyle{remark}
\newtheorem{remark}[theorem]{Remark}

\newcommand{\N}{\mathbb N}
\newcommand{\Nzero}{\mathbb N_0}
\newcommand{\Z}{\mathbb Z}
\newcommand{\Q}{\mathbb Q}
\newcommand{\R}{\mathbb R}
\newcommand{\C}{\mathbb C}
\newcommand{\E}{\mathcal E}
\newcommand{\cB}{\mathcal B}
\newcommand{\cC}{\mathcal C}
\newcommand{\cL}{\mathcal L}
\newcommand{\cP}{\mathcal P}
\newcommand{\cS}{\mathcal S}
\newcommand{\cV}{\mathcal V}
\newcommand{\Span}{\operatorname{span}}
\newcommand{\rank}{\operatorname{rank}}
\newcommand{\supp}{\operatorname{supp}}
\newcommand{\vol}{\operatorname{vol}}
\newcommand{\fract}[1]{\left\{#1\right\}}
\newcommand{\thetaAE}{\vartheta}
\newcommand{\statusproved}{\textcolor{provedgreen}{\textbf{[proved here]}}}
\newcommand{\statusconditional}{\textcolor{conditionalblue}{\textbf{[conditional on a counterexample]}}}
\newcommand{\statusopen}{\textcolor{opengray}{\textbf{[open target]}}}
\newcommand{\statuscompute}{\textcolor{computeviolet}{\textbf{[proposed computation]}}}
\newcommand{\statusimported}{\textcolor{conditionalblue}{\textbf{[imported from Report IV]}}}

\newenvironment{statusbox}[1]
{\begin{center}\begin{minipage}{0.94\textwidth}\raggedright\hrule\vspace{0.45em}\textbf{Status:} #1\par\vspace{0.35em}}
{\vspace{0.35em}\hrule\end{minipage}\end{center}}

\lstdefinestyle{pseudocode}{
  basicstyle=\ttfamily\footnotesize,
  backgroundcolor=\color{codegray},
  frame=single,
  rulecolor=\color{rulegray},
  showstringspaces=false,
  breaklines=true,
  columns=fullflexible,
  keepspaces=true,
  tabsize=2
}

\begin{document}
\pagenumbering{gobble}

\begin{titlepage}
\centering
\vspace*{0.6cm}
{\Huge\bfseries Algebraic Traps on the Power Curve\par}
\vspace{0.35cm}
{\LARGE New Progress Toward the Alaoglu--Erd\H{o}s Conjecture\par}
\vspace{0.8cm}
{\Large Deterministic Coppersmith--Chebyshev Certificates,\par}
\vspace{0.12cm}
{\Large Sparse Total Positivity, and the High-Corank Frontier\par}
\vspace{1.0cm}
{\large Progress report V --- 23 August 2026\par}
\vfill
\begin{minipage}{0.92\textwidth}
\small
\textbf{Bottom line.}
This report does not claim a proof of the Alaoglu--Erd\H{o}s conjecture.  It develops a new
and fully quantitative route prompted by the June 2026 work of Brisebarre and Hanrot on
Coppersmith-style detection of integer points near transcendental curves.  A hypothetical
counterexample supplies, in every logarithmic unit block, an initial segment of an irrational
rotation orbit on the fixed arc $y=x^{\log_2 3}$, sampled on the rational grid with denominators
$2^H$ and $3^H$.  On this particular curve, the usual ``two auxiliary polynomials must be
coprime'' heuristic can be removed: restrictions of sparse bivariate polynomials are
exponential sums with distinct real frequencies, hence form a strictly totally positive
Chebyshev system.  If $r$ linearly independent auxiliary polynomials with a common $q$-monomial
support are uniformly smaller than one, then the block contains at most $q-r$ exact integer
points.

A one-dimensional Chebyshev lattice, optimized over the $q$ lowest frequencies
$i\log 2+j\log 3$, then gives the unconditional block estimate
\[
 \#\{x\in[H,H+1):2^x,3^x\in\Z\}
 \le \left(\frac{32\log2\log3}{9}+o(1)\right)
       \frac{H^2}{(\log H)^2}.
\]
The logarithmic-square saving over a quadratic block bound is genuine new progress, but it
still does not contradict the linear block population forced by a counterexample.  Optimizing
macroscopic windows makes the residual barrier exact: with the current kernel-level lower
bound $\beta\ge14$ for a least counterexample, a closing trap would need to certify more than
$98.59\%$ of the entire coefficient space as uniformly short.  The remaining problem is thus
not ``find one clever auxiliary polynomial'' but ``prove a near-full-rank collapse of a
structured diagonal-flow lattice.''  The report formulates this as a parametric-geometry-of-
numbers problem, proves a dilation-rigidity theorem that constrains multiscale reuse, and gives
a proof-producing finite-certificate design suitable for exact arithmetic and Lean replay.
\end{minipage}
\vfill
{\small All new assertions are explicitly tagged.  No numerical reconnaissance is promoted to a theorem.\par}
\end{titlepage}

\clearpage
\pagenumbering{roman}
\begin{abstract}
Let
\[
 I(H):=\#\{x\in[H,H+1):2^x\in\Z,\ 3^x\in\Z\}.
\]
The attached unified report already develops the standard transcendence reductions, the exact
conditional solution monoid, determinant and $p$-adic methods, Kummer and cyclotomic
reformulations, compact-orbit dynamics, rational-point counting, and a large collection of
quantified no-go theorems.  The present document deliberately imports rather than repeats that
material.

Its first new observation is algorithmic: after rescaling the $H$th block by $(2^H,3^H)$,
all exact points lie on one fixed analytic arc and fit the rational-grid problem studied by
Brisebarre--Hanrot.  Its second observation is structural: for any finite monomial support
$E\subset\Nzero^2$, the functions
$e^{(i\log2+j\log3)t}$, $(i,j)\in E$, form an extended complete Chebyshev system.  This turns
short auxiliary vectors directly into deterministic point-counting certificates and replaces
coprimality by linear algebra.  Its third observation is asymptotic: a one-variable Chebyshev
expansion along the exact curve, combined with Cauchy--Binet and Minkowski's second theorem,
produces short integer polynomials.  The asymptotically optimal support is the weighted simplex
of the lowest frequencies, not the usual total-degree triangle.  This yields
\[
 \limsup_{H\to\infty} I(H)\frac{(\log H)^2}{H^2}
 \le \frac{32\log2\log3}{9}=2.7075555926\ldots.
\]

For a window $[H,(1+r)H]$, the same calculation has a sharp variational constant.  Writing
$t_0$ for the positive root of $t\tanh t=2$, the best coefficient-space-to-counterexample
population ratio is
\[
 \kappa_*=\frac{64\log2\log3}{9}\frac{\cosh t_0}{t_0^2}
 =5.0872665807\ldots.
\]
Thus a least counterexample $\beta\ge14$ would require certified short-vector rank fraction
strictly exceeding $1-1/(14\kappa_*)=0.9859593\ldots$.  This high-corank requirement is the
principal new frontier.  We connect it to Schmidt--Summerer--Roy parametric geometry of numbers,
prove exact dilation rigidity for sparse trap spaces, analyze why naive Hermite multiplicity is
not free, and specify a deterministic certificate format whose verification uses only integer
linear algebra and interval bounds.
\end{abstract}

\tableofcontents
\clearpage
\pagenumbering{arabic}

\section{Executive summary and status ledger}

\subsection{What is new in this report}

The mathematical contributions are the following.

\begin{enumerate}
\item \textbf{Exact curve-grid interface.}  Each logarithmic block is identified with exact
integer points on a fixed power arc sampled at the rational grid $(2^{-H}\Z)\times(3^{-H}\Z)$.
This is the precise interface with the 2026 Coppersmith--Chebyshev method.

\item \textbf{Sparse Chebyshev rigidity.}  A $q$-monomial polynomial restricts to an
exponential sum with $q$ distinct real frequencies and therefore has at most $q-1$ zeros on
the positive power curve, counted with multiplicity.

\item \textbf{Rank--codimension certificates.}  If $r$ independent integer polynomials with
one $q$-monomial support are uniformly smaller than one on an interval of the curve, then that
interval contains at most $q-r$ exact integer points.  No resultant and no coprimality
assumption is required.

\item \textbf{An optimized unconditional block theorem.}  A one-dimensional Chebyshev lattice
and a weighted-frequency simplex give
\[
 I(H)\le\left(\frac{32\log2\log3}{9}+o(1)\right)\frac{H^2}{(\log H)^2}.
\]
This is a paper proof, not a numerical heuristic.

\item \textbf{A sharp high-corank threshold.}  The best macroscopic-window constants are
solved explicitly.  The current finite exclusion $\beta\ge14$ forces any closing certificate
of this type to find at least $98.59\%$ of all coefficient directions below the shortness
threshold.

\item \textbf{A multiscale rigidity package.}  Trap spaces cannot recur exactly under the
$(X,Y)\mapsto(2X,3Y)$ dilation unless they are monomial.  This identifies the precise amount
of approximate covariance that a parametric-geometry argument would need.

\item \textbf{A deterministic certificate design.}  The report gives an exact output format
for interval-arithmetic Chebyshev bounds, integer coefficient vectors, and rank checks.  Such a
certificate proves a finite block bound without trusting the LLL implementation that found it.
\end{enumerate}

\subsection{What is not claimed}

No theorem below proves that a counterexample cannot exist.  In particular:

\begin{itemize}
\item the new $H^2/(\log H)^2$ block bound is still much larger than the conditional
$H/\beta+O(1)$ population;
\item no near-full-rank successive-minimum theorem is proved;
\item the proposed finite verifier has not been used here to enlarge the kernel-verified
zero-free region;
\item the 2026 AI breakthroughs discussed in the public record are motivational context only
and are not mathematical inputs.
\end{itemize}

\subsection{Status convention}

\begin{longtable}{p{0.24\textwidth}p{0.68\textwidth}}
\toprule
Tag & Meaning \\
\midrule
\endfirsthead
\toprule
Tag & Meaning \\
\midrule
\endhead
\statusproved & A complete conventional proof is given in this report. \\
\statusimported & Used from the attached unified report and not reproved here. \\
\statusconditional & Deduced under the assumption that a nonintegral solution exists. \\
\statuscompute & A concrete exact computation or implementation is proposed, but no output is
asserted as a theorem. \\
\statusopen & A precisely stated missing theorem or quantitative target. \\
\bottomrule
\end{longtable}

\section{Scope, nonduplication, and the 2026 trigger}

\subsection{The inherited frontier}

The attached Report IV is approximately twenty-four thousand lines of LaTeX.  It already
contains the elementary and transcendence reductions, kernel-verified finite exclusions, the
canonical primitive generator, exact block counts, radical and Kummer reformulations, auxiliary
base spectra, finite differences, generalized Vandermonde determinants, $p$-adic and tropical
ceilings, compact-orbit dynamics, o-minimal counting, Loewner and Pad\'e certificates, Mahler
systems, cyclotomic and Frobenius constructions, factorial cocycles, moment masks, Bhargava
factorials, and a detailed formalization inventory.  Repeating any of those developments would
make this report longer without moving the frontier.

Accordingly, only one interface theorem is imported: if the conjecture fails, all solutions
form the free additive monoid generated by $1$ and a unique least noninteger solution $\beta$.
Everything after Section~\ref{sec:imported-interface} is organized around a method absent from
Report IV: exact algebraic trapping of the power curve by short vectors in a Chebyshev lattice.

\subsection{Why the June 2026 paper matters}

Brisebarre and Hanrot introduced an algorithmic Coppersmith-style method for locating integer
points near a transcendental analytic curve, with Chebyshev interpolation controlling the
analytic error and lattice reduction producing auxiliary integer polynomials
\cite{BrisebarreHanrot2026}.  Their general algorithm is heuristic at one explicit point: two
short auxiliary polynomials can have a common factor, so a resultant computation may fail to
isolate the points.

The Alaoglu--Erd\H{o}s curve has extra structure that removes that difficulty.  A polynomial
restricted to $(2^t,3^t)$ becomes an exponential polynomial with pairwise distinct real
frequencies.  Such functions are strictly totally positive, so a single nonzero auxiliary
polynomial has a sharp zero bound, and several independent polynomials give an exact
rank--codimension bound.  The general Coppersmith machinery is therefore not merely applicable;
it simplifies in a way specific to multiplicatively independent bases.

A second 2026 input, due to Wang and Xiao, develops degeneracy theorems for families of unit
equations via integral points and new gcd estimates for exponential Diophantine problems
\cite{WangXiao2026}.  We use it only to motivate the stable-factor program of
Section~\ref{sec:stable-factors}; none of the unconditional results depends on it.

\subsection{Competitive context is not evidence}

The public 2026 record does substantiate striking AI-assisted mathematics, including an
explicit three-dimensional counterexample to the Jacobian conjecture discussed in detail by
Tao, and reported constructions covering previously unresolved Hadamard orders below $2000$
\cite{TaoJacobian2026,EpochHadamard2026}.  Those events justify urgency, not lowered standards.
Every claimed advance in this report is proved independently of model attribution, and every
remaining gap is stated rather than hidden.

\section{The imported conditional interface}
\label{sec:imported-interface}

Set
\[
 \cS:=\{x\in\R:2^x\in\N,\ 3^x\in\N\}.
\]

\begin{conjecture}[Alaoglu--Erd\H{o}s]
\label{conj:AE}
If $x\in\R$ and $2^x,3^x\in\Z$, then $x\in\Z$.
\end{conjecture}

The report uses the following exact conditional package.

\begin{theorem}[Primitive solution monoid; imported]
\label{thm:primitive-monoid}
\statusimported\quad
If Conjecture~\ref{conj:AE} is false, there is a unique least noninteger
$\beta\in\cS$ and integers
\[
 M=2^\beta,\qquad A=3^\beta
\]
such that
\[
 \cS=\{n+k\beta:n,k\in\Nzero\},
\]
and every representation is unique.  Consequently
\[
 (2^{n+k\beta},3^{n+k\beta})=(2^nM^k,3^nA^k).
\]
The attached report also gives the kernel-level finite exclusion $\beta\ge14$ and a
software-level exclusion $\beta>27$.
\end{theorem}

We shall use the exact unit-block count, also imported:
\begin{equation}
 \#(\cS\cap[H,H+1))=1+\left\lfloor\frac{H+1}{\beta}\right\rfloor
 \qquad(H\in\Nzero).
 \label{eq:exact-block-count}
\end{equation}
The $1$ is the integral solution $H$.

\subsection{Normalized block points}

Fix an integer $H\ge0$ and define
\[
 K_H:=\left\lfloor\frac{H+1}{\beta}\right\rfloor.
\]
For $1\le k\le K_H$, the unique solution in the block with second coordinate $k$ is
\[
 z_{H,k}=\lceil H-k\beta\rceil+k\beta=H+\fract{k\beta}.
\]
Its integer output pair is
\begin{align}
 X_{H,k}&=2^{z_{H,k}}=2^{H-\lfloor k\beta\rfloor}M^k,\label{eq:XHk}\\
 Y_{H,k}&=3^{z_{H,k}}=3^{H-\lfloor k\beta\rfloor}A^k.\label{eq:YHk}
\end{align}
After dividing by the block scales,
\begin{equation}
 \left(\frac{X_{H,k}}{2^H},\frac{Y_{H,k}}{3^H}\right)
 =\left(2^{\fract{k\beta}},3^{\fract{k\beta}}\right).
 \label{eq:normalized-orbit}
\end{equation}
Thus the normalized block is the initial orbit segment
\begin{equation}
 \cC_H=\{(1,1)\}\cup
 \left\{\left(2^{\fract{k\beta}},3^{\fract{k\beta}}\right):1\le k\le K_H\right\}
 \subset\{(2^t,3^t):0\le t<1\}.
 \label{eq:CH}
\end{equation}
The sets are nested: $\cC_H\subseteq\cC_{H+1}$, and at most one new nonintegral point is
inserted at each step.

\subsection{The exact rational-grid formulation}

Put
\[
 \thetaAE:=\frac{\log3}{\log2}.
\]
For $X=2^{H+t}$ and $Y=3^{H+t}$ one has
\[
 \frac{Y}{3^H}=\left(\frac{X}{2^H}\right)^{\thetaAE}.
\]
Therefore every exact point in the $H$th block is an integer pair $(X,Y)$ satisfying
\begin{equation}
 2^H\le X<2^{H+1},\qquad
 \frac{Y}{3^H}=f\!\left(\frac{X}{2^H}\right),
 \qquad f(u)=u^{\thetaAE},\quad 1\le u<2.
 \label{eq:grid-formulation}
\end{equation}
This is the Brisebarre--Hanrot rational-grid problem with fixed denominators
$u=2^H$, $v=3^H$ and zero vertical error.  Exactness is crucial: an auxiliary polynomial
$P\in\Z[X,Y]$ satisfying
\[
 \sup_{0\le t\le1}|P(2^{H+t},3^{H+t})|<1
\]
must vanish at every block point because its value there is an integer.

\begin{statusbox}{\statusconditional}
Equations~\eqref{eq:XHk}--\eqref{eq:CH} use the imported primitive-generator theorem.  The
rational-grid identity~\eqref{eq:grid-formulation} is unconditional and applies to every exact
integer point of the power curve, whether or not a primitive generator exists.
\end{statusbox}

\section{Sparse total positivity on the power curve}
\label{sec:total-positivity}

Let $E\subset\Nzero^2$ be finite.  Write
\[
 q:=|E|,\qquad \lambda_{i,j}:=i\log2+j\log3.
\]
The frequencies $\lambda_{i,j}$ are pairwise distinct: equality of two frequencies would give
$2^{i-i'}3^{j-j'}=1$, hence $i=i'$ and $j=j'$.

\subsection{The exponential Chebyshev system}

\begin{theorem}[Strict total positivity]
\label{thm:strict-total-positivity}
Let
\[
 \lambda_1<\cdots<\lambda_q,
 \qquad t_1<\cdots<t_m.
\]
The matrix
\[
 V=(e^{\lambda_jt_i})_{1\le i\le m,\ 1\le j\le q}
\]
has rank $\min(m,q)$.  More strongly, every minor obtained from increasing row and column
indices is strictly positive.
\end{theorem}

\begin{proof}
It is enough to prove that every square generalized Vandermonde determinant
\[
 \det(e^{\mu_j s_i})_{1\le i,j\le r}
\]
is positive for $s_1<\cdots<s_r$ and $\mu_1<\cdots<\mu_r$.  Divide column $j$ by
$e^{\mu_1s_j}$ after transposing if desired; equivalently, use the classical fact that the
Wronskian of $e^{\mu_1t},\ldots,e^{\mu_rt}$ is
\[
 e^{(\mu_1+\cdots+\mu_r)t}\prod_{1\le i<j\le r}(\mu_j-\mu_i)>0.
\]
A family of $C^{r-1}$ functions whose initial Wronskians are everywhere positive is an
extended complete Chebyshev system.  Its collocation determinants at increasing nodes are
positive.  For completeness, this last implication can be proved inductively: divide a
nontrivial linear combination by the first exponential, differentiate, and apply Rolle's
theorem.  It follows that no nontrivial combination of $r$ functions can have $r$ real zeros
counted with multiplicity; expansion of a collocation determinant along a hypothetical linear
dependence would contradict that zero bound.  The sign is fixed by coalescing the nodes, where
the determinant tends to the positive Wronskian times the positive confluent Vandermonde
factor.
\end{proof}

\begin{corollary}[Sparse zero bound]
\label{cor:sparse-zero-bound}
Let
\[
 P(X,Y)=\sum_{(i,j)\in E}c_{i,j}X^iY^j\in\R[X,Y]
\]
be nonzero.  Then
\[
 F_P(t):=P(2^t,3^t)
       =\sum_{(i,j)\in E}c_{i,j}e^{\lambda_{i,j}t}
\]
has at most $q-1$ real zeros, counted with multiplicity.
\end{corollary}

\begin{proof}
Order the distinct frequencies and apply the extended Chebyshev property from
Theorem~\ref{thm:strict-total-positivity}.
\end{proof}

\begin{remark}[Why this is stronger than B\'ezout here]
B\'ezout controls common zeros of two coprime algebraic curves.  Corollary~\ref{cor:sparse-zero-bound}
controls the intersection of a \emph{single} sparse polynomial with the transcendental power
curve and remains valid even if several auxiliary polynomials share factors.  The relevant
complexity is the number of occupied monomials, not the total algebraic degree.
\end{remark}

\subsection{Multiplicity and evaluation rank}

For distinct real nodes $t_1,\ldots,t_m$, let
\[
 \operatorname{ev}_{\mathbf t}:\R^E\longrightarrow\R^m,
 \qquad
 (c_e)_{e\in E}\longmapsto
 \left(\sum_{e\in E}c_e e^{\lambda_e t_a}\right)_{a=1}^m.
\]

\begin{corollary}[Exact evaluation rank]
\label{cor:evaluation-rank}
For distinct nodes, $\rank(\operatorname{ev}_{\mathbf t})=\min(m,q)$.  Hence, when $m<q$,
\[
 \dim\ker(\operatorname{ev}_{\mathbf t})=q-m.
\]
\end{corollary}

This elementary dimension identity is the engine behind the deterministic trap theorem.

\section{Rank--codimension algebraic traps}
\label{sec:rank-codim}

\subsection{A deterministic replacement for coprimality}

\begin{theorem}[Rank--codimension trap]
\label{thm:rank-codim}
Fix a finite support $E\subset\Nzero^2$ of size $q$, a real interval $J$, and a scale $H$.
Suppose
\[
 P_1,\ldots,P_r\in\Z[X,Y]
\]
have coefficient vectors supported in $E$, those vectors are linearly independent over
$\R$, and
\begin{equation}
 \sup_{t\in J}|P_a(2^{H+t},3^{H+t})|<1
 \qquad(1\le a\le r).
 \label{eq:trap-shortness}
\end{equation}
Then the curve segment
\[
 \{(2^{H+t},3^{H+t}):t\in J\}
\]
contains at most $q-r$ points of $\Z^2$.
\end{theorem}

\begin{proof}
Let $t_1,\ldots,t_m\in J$ be the parameters of distinct integer points.  Since each
$P_a(2^{H+t_b},3^{H+t_b})$ is an integer of absolute value less than one, it is zero.  Thus the
$r$ coefficient vectors lie in the kernel of the evaluation map at the $m$ nodes.  By
Corollary~\ref{cor:evaluation-rank}, this kernel has dimension $q-m$ (and is zero if $m\ge q$).
Therefore $r\le q-m$, or $m\le q-r$.
\end{proof}

\begin{statusbox}{\statusproved}
Theorem~\ref{thm:rank-codim} is exact and deterministic.  It needs neither a resultant nor a
coprimality test.  The only numerical statement to certify is the strict bound
\eqref{eq:trap-shortness}; the rank of the integer coefficient vectors is checked exactly.
\end{statusbox}

\begin{corollary}[One polynomial]
\label{cor:one-polynomial}
A single nonzero $q$-monomial polynomial satisfying~\eqref{eq:trap-shortness} bounds the number
of exact integer points by $q-1$.
\end{corollary}

\begin{corollary}[Several supports]
\label{cor:union-support}
Suppose $P_1,\ldots,P_r$ have possibly different supports.  Let $E$ be their union and let
$r'$ be the rank of their coefficient vectors after padding by zero in $\R^E$.  Then the same
segment contains at most $|E|-r'$ integer points.
\end{corollary}

\subsection{Piecewise certificates}

Global uniform approximation is not mandatory.  Let $J_1,\ldots,J_s$ be intervals with
disjoint interiors covering $J$.  On $J_\nu$, let $E_\nu$ be a support of size $q_\nu$ and
let $r_\nu$ be the rank of certified short polynomials.

\begin{criterion}[Cover certificate]
\label{crit:cover}
The number of exact integer points with parameter in $J$ is at most
\begin{equation}
 \sum_{\nu=1}^s(q_\nu-r_\nu),
 \label{eq:cover-defect}
\end{equation}
with endpoints assigned to one neighboring interval to avoid double counting.
\end{criterion}

The quantity $q_\nu-r_\nu$ is the \emph{local coefficient defect}.  Formula
\eqref{eq:cover-defect} turns the search for points into a search for low total defect.

\begin{corollary}[Conditional closing criterion]
\label{cor:conditional-closing}
Assume a counterexample exists with primitive generator $\beta$.  If, for infinitely many
integers $H$, one can cover $[0,1]$ by certified traps whose total defect is
$o(H)$, then Conjecture~\ref{conj:AE} follows.
\end{corollary}

\begin{proof}
The exact block count~\eqref{eq:exact-block-count} is $H/\beta+O(1)$, while the cover
certificate gives $o(H)$.
\end{proof}

\subsection{Successive minima are the correct invariant}

Every analytic construction below associates to $(H,J,E)$ a lattice $\cL_{H,J,E}$ whose
short vectors are coefficient vectors of uniformly small integer polynomials.  Let
\[
 0<\mu_1\le\cdots\le\mu_q
\]
be its successive minima relative to the norm that dominates the curve sup norm, and let
$\tau$ be the strict threshold corresponding to ``supremum $<1$.''

\begin{criterion}[Successive-minimum point bound]
\label{crit:successive-minima}
If $\mu_r<\tau$, then the interval contains at most $q-r$ exact integer points.  Conversely,
if the interval contains $m$ exact integer points, then
\[
 \mu_{q-m+1}\ge\tau.
\]
\end{criterion}

Thus the problem is spectral.  Finding the first short vector proves a $q-1$ bound; closing the
conjecture requires control near the \emph{top} of the successive-minimum spectrum.

\section{A one-dimensional Chebyshev lattice on the exact curve}
\label{sec:cheb-lattice}

The general Brisebarre--Hanrot construction approximates a two-dimensional strip around a
curve.  Exact equality permits a leaner object: expand each monomial directly along
\[
 \Phi_H(t)=(2^{H+t},3^{H+t}),\qquad 0\le t\le1.
\]
This reduces the analytic part to one variable and exposes the frequency geometry.

\subsection{Bernstein ellipses and coefficient bounds}

Map $t\in[0,1]$ to $u=2t-1\in[-1,1]$.  For $\rho>1$, let $\mathscr E_\rho$ be the Bernstein
ellipse in the $u$-plane and put
\begin{equation}
 R(\rho):=\frac12+\frac{\rho+\rho^{-1}}4.
 \label{eq:Rrho}
\end{equation}
The maximum real part of $t=(1+u)/2$ on $\mathscr E_\rho$ is $R(\rho)$.

For $e=(i,j)\in E$ define
\[
 \lambda_e=i\log2+j\log3,
 \qquad
 f_e(t)=e^{\lambda_e(H+t)}.
\]
Write its Chebyshev expansion
\[
 f_e(t)=\sum_{n\ge0}' a_{n,e}T_n(2t-1),
\]
where the prime only changes the conventional normalization of the constant coefficient.
Standard contour estimates give, after harmless adjustment of the $n=0$ factor,
\begin{equation}
 |a_{n,e}|\le 2M_e\rho^{-n},
 \qquad
 M_e:=e^{\lambda_e(H+R(\rho))}.
 \label{eq:cheb-coeff-bound}
\end{equation}
Consequently the tail after $q$ terms is bounded by
\begin{equation}
 \sum_{n\ge q}|a_{n,e}|
 \le \tau_e,
 \qquad
 \tau_e:=\frac{2M_e\rho^{-q}}{1-\rho^{-1}}.
 \label{eq:tail-weight}
\end{equation}

\subsection{The coefficient-and-tail matrix}

Let $A$ be the $q\times q$ real matrix with entries $A_{n,e}=a_{n,e}$,
$0\le n<q$, and let $D=\operatorname{diag}(\tau_e)$.  Define the $2q\times q$ matrix
\begin{equation}
 B=B(H,E,\rho):=\begin{pmatrix}A\\D\end{pmatrix}
 \label{eq:Bmatrix}
\end{equation}
and the rank-$q$ lattice
\[
 \cL(H,E,\rho):=B\Z^E\subset\R^{2q}.
\]
For $c=(c_e)\in\Z^E$, let
\[
 P_c(X,Y)=\sum_{e=(i,j)\in E}c_eX^iY^j.
\]

\begin{lemma}[Norm-to-supremum transfer]
\label{lem:norm-to-sup}
For every $c\in\Z^E$,
\begin{equation}
 \sup_{0\le t\le1}|P_c(2^{H+t},3^{H+t})|
 \le \sqrt{2q}\,\|Bc\|_2.
 \label{eq:norm-to-sup}
\end{equation}
\end{lemma}

\begin{proof}
The first $q$ Chebyshev coefficients of the linear combination are $Ac$.  Their absolute sum
is at most $\sqrt q\|Ac\|_2$.  By~\eqref{eq:tail-weight}, the tail is at most
$\sum_e|c_e|\tau_e\le\sqrt q\|Dc\|_2$.  Since $|T_n(u)|\le1$ on $[-1,1]$, the sum of these
bounds is at most $\sqrt{2q}\|Bc\|_2$.
\end{proof}

Thus a lattice vector shorter than $(2q)^{-1/2}$ is a valid algebraic trap.

\subsection{Covolume by Cauchy--Binet}

Let
\[
 \Delta(H,E,\rho):=\sqrt{\det(B^TB)}
\]
be the covolume of $\cL(H,E,\rho)$ in its span.

\begin{lemma}[Universal covolume bound]
\label{lem:covolume}
Assume $\rho\ge2$.  Uniformly over all finite supports $E$ of size $q$,
\begin{equation}
 \log\Delta(H,E,\rho)
 \le (H+R(\rho))\sum_{e\in E}\lambda_e
      -\frac{q(q-1)}2\log\rho+O(q\log q),
 \label{eq:covolume-bound}
\end{equation}
where the implied constant is absolute.
\end{lemma}

\begin{proof}
Cauchy--Binet writes $\det(B^TB)$ as the sum of squares of all $q\times q$ minors of $B$.
Choose $k$ rows from $A$ and $q-k$ diagonal rows from $D$.  A nonzero minor factors into a
$k\times k$ coefficient minor and $q-k$ diagonal weights.  By
\eqref{eq:cheb-coeff-bound}, Hadamard's inequality bounds its absolute value by
\[
 e^{(H+R)\sum_{e\in E}\lambda_e}\,
 \rho^{-\sum_{n\in S}n-q(q-k)}\,e^{O(q\log q)},
\]
where $S$ is the chosen $k$-set of Chebyshev indices.  The smallest possible exponent of
$\rho^{-1}$ is obtained from $S=\{0,\ldots,k-1\}$, and
\begin{align*}
 q(q-k)+\frac{k(k-1)}2
 &=\frac{q(q-1)}2+\frac{(q-k)(q-k+1)}2\\
 &\ge\frac{q(q-1)}2.
\end{align*}
The number of minors and all factors $2$, $(1-\rho^{-1})^{-1}$, and Hadamard constants
contribute $e^{O(q\log q)}$.  Taking the square root of the Cauchy--Binet sum gives
\eqref{eq:covolume-bound}.
\end{proof}

The identity in the proof is important: diagonal tail rows never reduce the full
$q(q-1)/2$ Chebyshev decay.  They only add further decay.

\subsection{Minkowski and fixed-rank short vectors}

Let $\mu_1\le\cdots\le\mu_q$ be the Euclidean successive minima of $\cL(H,E,\rho)$.
Minkowski's second theorem gives
\[
 \mu_1\cdots\mu_q\,\vol(B_2^q)\le2^q\Delta(H,E,\rho).
\]
In particular,
\begin{equation}
 \log\mu_1
 \le \frac{H+R(\rho)}q\sum_{e\in E}\lambda_e
      -\frac{q-1}{2}\log\rho+O(\log q).
 \label{eq:first-minimum}
\end{equation}

To obtain several vectors, note that every nonzero $c\in\Z^E$ obeys
\[
 \|Bc\|_2\ge\min_e\tau_e.
\]
When $(0,0)\in E$, the minimum is at least $2\rho^{-q}$ for $\rho\ge2$.  Therefore, for each
fixed $r$,
\begin{equation}
 \log\mu_r
 \le \frac{H+R(\rho)}q\sum_{e\in E}\lambda_e
      -\frac q2\log\rho+O_r(\log q+\log\rho).
 \label{eq:fixed-r-minimum}
\end{equation}
The difference between $q/2$ and $(q-1)/2$ is absorbed in the error.

\begin{theorem}[Fixed-rank Chebyshev trap]
\label{thm:fixed-rank-trap}
Let $r$ be fixed.  Suppose $(0,0)\in E_H$, $q_H=|E_H|\to\infty$, and there are parameters
$\rho_H\ge2$ such that
\begin{equation}
 \frac{H+R(\rho_H)}{q_H}\sum_{e\in E_H}\lambda_e
 -\frac{q_H}{2}\log\rho_H\longrightarrow-\infty
 \label{eq:shortness-criterion}
\end{equation}
faster than $\log q_H+\log\rho_H$.  Then, for all sufficiently large $H$, there are $r$
linearly independent polynomials in $\Z[X,Y]$, supported in $E_H$, whose absolute values on
$\Phi_H([0,1])$ are strictly less than one.
\end{theorem}

\begin{proof}
Equation~\eqref{eq:fixed-r-minimum} makes $\mu_r<(2q_H)^{-1/2}$ for large $H$.  Choose $r$
independent lattice vectors below that threshold and apply Lemma~\ref{lem:norm-to-sup}.
\end{proof}

\begin{statusbox}{\statusproved}
Theorem~\ref{thm:fixed-rank-trap} is an existence theorem in the exact real Chebyshev lattice.
A proof-producing implementation replaces the real coefficients by outward-rounded rational
intervals and verifies the final sup-norm inequality independently; see
Section~\ref{sec:certificate-design}.
\end{statusbox}

\section{Frequency-simplex optimization and a new block bound}
\label{sec:block-bound}

The choice of monomial support determines both sides of
\eqref{eq:shortness-criterion}.  For fixed cardinality $q$, the negative Chebyshev term depends
only on $q$, while the positive growth term is
$\sum_{e\in E}\lambda_e$.  The optimal support is therefore obtained by taking the $q$
smallest frequencies.

Put
\[
 a:=\log2,
 \qquad b:=\log3.
\]
For $L>0$, define the weighted frequency simplex
\begin{equation}
 E_L:=\{(i,j)\in\Nzero^2:ai+bj\le L\}.
 \label{eq:weighted-simplex}
\end{equation}
Let
\[
 Q(L):=|E_L|,
 \qquad
 S(L):=\sum_{(i,j)\in E_L}(ai+bj).
\]

\subsection{Lattice-point asymptotics}

\begin{lemma}[Weighted simplex count and first moment]
\label{lem:simplex-asymptotics}
As $L\to\infty$,
\begin{align}
 Q(L)&=\frac{L^2}{2ab}+O(L),\label{eq:QL}\\
 S(L)&=\frac{L^3}{3ab}+O(L^2),\label{eq:SL}\\
 \frac{S(L)}{Q(L)}&=\frac{2L}{3}+O(1).\label{eq:average-frequency}
\end{align}
\end{lemma}

\begin{proof}
Summing vertical fibers gives
\[
 Q(L)=\sum_{0\le i\le L/a}
 \left(1+\left\lfloor\frac{L-ai}{b}\right\rfloor\right).
\]
Replacing each floor by its argument incurs $O(L)$, and the remaining arithmetic progression
is the area $L^2/(2ab)$ plus $O(L)$.  For the first moment, sum
\[
 \sum_i\sum_{0\le j\le(L-ai)/b}(ai+bj).
\]
Replacing the inner sum by its integral incurs $O(L^2)$.  The integral over the triangle
$ai+bj\le L$ is
\[
 \int_0^{L/a}\int_0^{(L-ai)/b}(ai+bj)\,dj\,di=\frac{L^3}{3ab}.
\]
Division yields~\eqref{eq:average-frequency}.
\end{proof}

\begin{proposition}[Asymptotic support optimality]
\label{prop:support-optimality}
Among all $q$-element subsets of $\Nzero^2$, the sum of frequencies is minimized by the
$q$ lattice points with smallest values of $ai+bj$.  Hence weighted simplices are
asymptotically optimal for the Chebyshev covolume problem.
\end{proposition}

\begin{proof}
This is the rearrangement principle: if a support contains a frequency larger than an omitted
one, swapping the two decreases the sum without changing $q$.
\end{proof}

The usual total-degree triangle $i+j\le d$ has average frequency
$d(a+b)/3$.  At the same cardinality, the weighted simplex has asymptotic average
$2\sqrt{2abq}/3$, improving the leading constant by the factor
\[
 \frac{2\sqrt{ab}}{a+b}=0.973\ldots.
\]
The gain is modest but exact and costs nothing.

\subsection{Choice of the Bernstein ellipse}

Insert $E_L$ into~\eqref{eq:first-minimum}.  By
Lemma~\ref{lem:simplex-asymptotics},
\begin{equation}
 \log\mu_1
 \le \frac{2L}{3}\bigl(H+R(\rho)\bigr)
      -\frac{L^2}{4ab}\log\rho
      +O(L^2+\log L).
 \label{eq:weighted-minimum}
\end{equation}
The $O(L^2)$ term absorbs the boundary errors from $Q(L)$ and $S(L)$.

For the unit block, the optimal ellipse grows with $L$.  Since
$R(\rho)=\rho/4+1/2+O(\rho^{-1})$, differentiating the two $\rho$-dependent leading terms
suggests
\begin{equation}
 \rho=\frac{3L}{2ab}.
 \label{eq:rho-optimal-unit}
\end{equation}
With this choice,
\begin{equation}
 \log\mu_1
 \le \frac{2LH}{3}-\frac{L^2}{4ab}\log L+O(L^2).
 \label{eq:unit-minimum-final}
\end{equation}
The same leading estimate holds for every fixed successive minimum $\mu_r$ by
\eqref{eq:fixed-r-minimum}.

\subsection{The logarithmic-square saving}

\begin{theorem}[Unconditional integer-point bound in a logarithmic unit block]
\label{thm:block-bound}
Let
\[
 I(H):=\#\{x\in[H,H+1):2^x\in\Z,\ 3^x\in\Z\}.
\]
Then
\begin{equation}
 \boxed{
 \limsup_{H\to\infty}
 I(H)\frac{(\log H)^2}{H^2}
 \le \frac{32\log2\log3}{9}}
 \label{eq:block-bound-main}
\end{equation}
and
\[
 \frac{32\log2\log3}{9}=2.70755559260021\ldots.
\]
More precisely, for every $\varepsilon>0$ and every fixed $r\ge1$, all sufficiently large
$H$ admit $r$ independent uniformly short integer polynomials supported on a set of size
\begin{equation}
 \left(\frac{32\log2\log3}{9}+\varepsilon+o(1)\right)
 \frac{H^2}{(\log H)^2},
 \label{eq:q-size}
\end{equation}
so Theorem~\ref{thm:rank-codim} subtracts $r$ from that bound.
\end{theorem}

\begin{proof}
Fix $\delta>0$ and choose
\begin{equation}
 L=\left(\frac{8ab}{3}+\delta\right)\frac{H}{\log H},
 \qquad
 \rho=\frac{3L}{2ab}.
 \label{eq:Lchoice}
\end{equation}
Since $\log L=\log H+O(\log\log H)$, the two leading terms in
\eqref{eq:unit-minimum-final} equal
\begin{align*}
 \frac{2LH}{3}-\frac{L^2}{4ab}\log L
 &=-c_\delta\frac{H^2}{\log H}
   +O\!\left(\frac{H^2\log\log H}{(\log H)^2}\right)
\end{align*}
for some $c_\delta>0$.  Thus the fixed-rank shortness criterion
\eqref{eq:shortness-criterion} holds with a margin tending to $-\infty$.  Theorem
\ref{thm:fixed-rank-trap} and then Theorem~\ref{thm:rank-codim} give
$I(H)\le Q(L)-r$.

Finally,
\[
 Q(L)=\frac{L^2}{2ab}+O(L)
 =\left(\frac{32ab}{9}+O(\delta)\right)
   \frac{H^2}{(\log H)^2}.
\]
Let $\delta\downarrow0$.
\end{proof}

\begin{statusbox}{\statusproved}
Theorem~\ref{thm:block-bound} is the principal unconditional advance of this report.  It is
not a reformulation of the conjecture and does not assume the primitive solution monoid.  Its
input is only multiplicative independence of $2$ and $3$, Chebyshev coefficient decay,
Cauchy--Binet, Minkowski's theorem, and strict total positivity.
\end{statusbox}

\subsection{Why the theorem still falls short}

Under a counterexample,~\eqref{eq:exact-block-count} gives
\[
 I(H)=\frac{H}{\beta}+O(1).
\]
The new upper bound is
\[
 I(H)\ll\frac{H^2}{(\log H)^2},
\]
which is smaller than a quadratic bound by a logarithmic square but still larger than the
conditional linear count by the factor $H/(\log H)^2$.  This is not a matter of constants.
A proof needs either:

\begin{itemize}
\item a much smaller coefficient support;
\item a rank $r$ so close to $q$ that the defect $q-r$ becomes $o(H)$; or
\item an additional arithmetic mechanism that forces multiplicity or shared factors without
paying the full coefficient-space dimension.
\end{itemize}

The next sections quantify the second possibility and audit the third.

\section{Macroscopic windows and the exact high-corank threshold}
\label{sec:window-optimization}

The unit-block theorem uses an ellipse whose radius grows like $H/\log H$.  To compare the
trap dimension directly with the quadratic cumulative population forced by a counterexample,
consider a macroscopic exponent window
\begin{equation}
 \mathcal I_{H,r}:=[H,(1+r)H],\qquad r>0.
 \label{eq:macro-window}
\end{equation}
Parameterize it by $z=H(1+rs)$, $0\le s\le1$.

\subsection{Trap dimension in a fixed-ratio window}

For a weighted simplex with $q$ monomials, its average frequency is
\begin{equation}
 \frac1q\sum_{e\in E}\lambda_e
 =\frac{2}{3}\sqrt{2abq}+o(\sqrt q).
 \label{eq:average-q}
\end{equation}
On a Bernstein ellipse of fixed parameter $\rho>1$, the largest real exponent is
$H(1+rR(\rho))$.  The first-minimum inequality becomes negative when
\begin{equation}
 \sqrt q>
 \frac{4\sqrt{2ab}}{3}
 \frac{1+rR(\rho)}{\log\rho}\,H+o(H).
 \label{eq:macro-q-threshold}
\end{equation}
Thus a fixed-rank trap can be built with
\begin{equation}
 q=\left(
 \frac{32ab}{9}\frac{(1+rR(\rho))^2}{(\log\rho)^2}+o(1)
 \right)H^2.
 \label{eq:macro-q}
\end{equation}

\subsection{Population supplied by a counterexample}

The cumulative conditional count below exponent $T$ is
$T^2/(2\beta)+O(T)$.  Therefore the window~\eqref{eq:macro-window} contains
\begin{equation}
 \left(\frac{r+r^2/2}{\beta}+o(1)\right)H^2
 \label{eq:macro-population}
\end{equation}
solutions.  The ratio between the fixed-rank trap dimension and this population is
asymptotically
\begin{equation}
 \beta\,\kappa(r,\rho),
 \qquad
 \kappa(r,\rho):=
 \frac{32ab}{9}
 \frac{(1+rR(\rho))^2}{(r+r^2/2)(\log\rho)^2}.
 \label{eq:kappa-r-rho}
\end{equation}

\subsection{Exact optimization}

\begin{theorem}[Optimal macroscopic-window constant]
\label{thm:kappa-opt}
Let $t_0>0$ be the unique solution of
\begin{equation}
 t\tanh t=2.
 \label{eq:t0}
\end{equation}
Then
\begin{equation}
 \kappa_*:=\inf_{r>0,\rho>1}\kappa(r,\rho)
 =\frac{64\log2\log3}{9}\frac{\cosh t_0}{t_0^2}
 =5.08726658072721\ldots.
 \label{eq:kappa-star}
\end{equation}
The minimizers are
\begin{align}
 t_0&=2.06533813897470\ldots,\label{eq:t0-num}\\
 \rho_0=e^{t_0}&=7.88796467504653\ldots,\label{eq:rho0}\\
 r_0&=0.665032892479158\ldots.
 \label{eq:r0}
\end{align}
\end{theorem}

\begin{proof}
For fixed $\rho$, write $R=R(\rho)$.  Differentiating
$(1+rR)^2/(r+r^2/2)$ gives the unique minimum
\[
 r=\frac1{R-1}.
\]
At this value,
\[
 \frac{(1+rR)^2}{r+r^2/2}=2(2R-1).
\]
Put $t=\log\rho$.  Since
\[
 2R(\rho)-1=\frac{\rho+\rho^{-1}}2=\cosh t,
\]
we obtain
\[
 \inf_r\kappa(r,e^t)=\frac{64ab}{9}\frac{\cosh t}{t^2}.
\]
The derivative vanishes exactly when $t\sinh t=2\cosh t$, or
$t\tanh t=2$.  The left side is strictly increasing from $0$ to infinity, so the root is
unique.  Substitution gives the numerical values.
\end{proof}

\subsection{The required rank fraction}

If a trap space has dimension $q$ and certified short rank $R_H$, its point bound is
$q-R_H$.  In the optimized macroscopic window, contradiction with
\eqref{eq:macro-population} requires
\begin{equation}
 \frac{R_H}{q}>1-\frac1{\kappa_*\beta}+o(1).
 \label{eq:rank-fraction}
\end{equation}

\begin{corollary}[Numerical high-corank frontier]
\label{cor:rank-frontier}
Using the imported kernel-level bound $\beta\ge14$, a closing trap must certify
\begin{equation}
 \frac{R_H}{q}>0.985959341761414\ldots,
 \label{eq:rank-kernel}
\end{equation}
that is, more than $98.59\%$ of all coefficient directions.  Using the existing
software-level exclusion $\beta>27$, the threshold becomes
\begin{equation}
 \frac{R_H}{q}>0.992719658691104\ldots.
 \label{eq:rank-software}
\end{equation}
\end{corollary}

\begin{statusbox}{\statusproved}
The constants in Theorem~\ref{thm:kappa-opt} and Corollary~\ref{cor:rank-frontier} are analytic,
not fitted numerically.  The only numerical operation is evaluation of the unique root of
$t\tanh t=2$.
\end{statusbox}

\subsection{Interpretation}

A two-polynomial Coppersmith calculation is nowhere near the required regime.  Even a theorem
producing ten, one hundred, or a fixed positive number of independent short vectors would not
change the asymptotic comparison.  The needed statement is qualitatively different: the
lattice must have only about $1.4\%$ of its directions above threshold at the kernel frontier,
and less than $0.73\%$ at the software frontier.  This is why the next natural language is
parametric geometry of numbers rather than ordinary LLL output.

\section{What determinant volume can and cannot say about rank}
\label{sec:volume-rank}

Minkowski's second theorem controls the product of all successive minima.  Turning that product
into a bound for a high individual minimum requires lower bounds for the earlier minima.  The
diagonal tail rows provide such a lower bound, but it is exponentially small and therefore
expensive when used repeatedly.

\subsection{A general rank-tax inequality}

Assume $(0,0)\in E$ and $\rho\ge2$.  From~\eqref{eq:tail-weight}, every nonzero lattice vector
has norm at least $2\rho^{-q}$.  If $1\le r\le q$, then
\[
 \mu_1\cdots\mu_q
 \ge (2\rho^{-q})^{r-1}\mu_r^{q-r+1}.
\]
Combining this with Minkowski's second theorem and Lemma~\ref{lem:covolume} yields
\begin{align}
 (q-r+1)\log\mu_r
 &\le (H+R(\rho))\sum_{e\in E}\lambda_e
      -\frac{q(q-1)}2\log\rho\notag\\
 &\quad +(r-1)q\log\rho+O(q\log q).
 \label{eq:rank-tax}
\end{align}
The term $(r-1)q\log\rho$ is the \emph{rank tax}.  It is absent for the first minimum and
lower order for fixed $r$, but it is quadratic in $q$ when $r$ is a positive fraction of the
dimension.

\subsection{The one-half ceiling of this proof mechanism}

Ignore the positive archimedean term in~\eqref{eq:rank-tax}; this only makes short vectors
easier.  If $r=\alpha q+o(q)$, the leading $\log\rho$ coefficient in the numerator is
\[
 \left(-\frac12+\alpha\right)q^2.
\]
Therefore this determinant-plus-universal-lower-bound argument cannot prove
$\mu_r<1$ for any fixed $\alpha>1/2$, even in the fictitious case where all monomial growth
vanishes.  This is a ceiling of the \emph{argument}, not a lower bound on the actual minima.
It shows that the $98.59\%$ target in~\eqref{eq:rank-kernel} cannot be reached by feeding the
same covolume estimate into Minkowski more aggressively.

\begin{warning}[Volume is not spectrum]
A small determinant can coexist with a few extremely short directions and many long ones.
Conversely, near-full short rank is a statement about the shape of the lattice, not merely its
covolume.  Any future claim that ``the determinant is tiny, hence almost every basis vector is
short'' must supply a genuine condition-number, transference, or successive-minimum theorem.
\end{warning}

\subsection{A conservation law in the unit-block asymptotic}

For the unit-block weighted simplex, let the outer frequency scale be
$L=cH/\log H$.  The determinant margin permits, through~\eqref{eq:rank-tax}, a rank fraction no
larger than approximately
\[
 \alpha(c)<\frac12-\frac{4ab}{3c},
 \qquad c>\frac{8ab}{3}.
\]
The residual coefficient defect has leading size
\[
 (1-\alpha(c))Q(L)
 \gtrsim
 \left(\frac12+\frac{4ab}{3c}\right)
 \frac{c^2}{2ab}\frac{H^2}{(\log H)^2}.
\]
The right side is minimized at the threshold $c=8ab/3$, where $\alpha=0$, and reproduces
exactly $32ab/9$.  Increasing the support makes more rank available, but the larger ambient
dimension cancels the gain.  Within this coarse volume method, fixed-rank and fractional-rank
optimization have the same leading point-count constant.

\section{Parametric geometry of the trap lattices}
\label{sec:pgn}

The scale $H$ enters each monomial column through
\[
 e^{\lambda_eH}.
\]
After separating the $H$-independent Chebyshev data, the family is a diagonal flow with weight
vector $(\lambda_e)_{e\in E}$.  This is the natural setting of Schmidt--Summerer parametric
geometry of numbers and Roy's rigid-system realization theory
\cite{SchmidtSummerer2009,SchmidtSummerer2013,Roy2015}.

\subsection{Successive-minimum functions}

For a fixed support and a fixed normalized interval, write
\[
 L_i(H):=\log\mu_i(H),\qquad 1\le i\le q.
\]
The columns of the defining lattice are multiplied by $e^{\lambda_eH}$ as $H$ varies.
Standard parametric-geometry heuristics predict piecewise-linear behavior for the ordered
functions $L_i(H)$, with slopes drawn from finite sums of the weights $\lambda_e$ and switch
points at which distinct short-vector supports exchange order.

Two features make this family more structured than a generic diagonal flow:

\begin{enumerate}
\item the weights are the two-dimensional semigroup $i\log2+j\log3$;
\item the analytic rows are Chebyshev coefficients of the same exponentials, so their
$H$-independent part is strictly totally positive.
\end{enumerate}

A high-corank theorem would follow from a strong upper bound on the upper tail of
$L_i(H)$, not from control of their sum alone.

\begin{target}[Trap-spectrum theorem]
\label{target:trap-spectrum}
For weighted-simplex supports in the optimized macroscopic window, prove one of the following:
\begin{enumerate}
\item at least $(1-1/(\kappa_*\beta)+\eta)q$ successive minima lie below the exact trap
threshold for infinitely many $H$;
\item the set of indices above threshold has cardinality $o(q)$;
\item a transference inequality bounds the long tail by arithmetic data at primes dividing
$6MA$ and forces the same conclusion under a hypothetical primitive pair $(M,A)$.
\end{enumerate}
Any one of these statements, with the first using $\beta\ge14$, proves the conjecture.
\end{target}

\subsection{Switch density as a finite combinatorial target}

The coefficient vector of a short polynomial has a support pattern, sign pattern, and a
leading Chebyshev cancellation pattern.  If only $o(q)$ directions are long, then almost all
patterns must persist over substantial $H$-intervals.  Conversely, if patterns switch too
often, pairs of nearly equal exponential weights arise.  Their gaps are integer linear forms
in $\log2$ and $\log3$:
\[
 |(i-i')\log2+(j-j')\log3|.
\]
Baker's theorem gives explicit lower bounds for nonzero gaps.  This suggests a concrete
program:

\begin{target}[Switch-count bound]
\label{target:switch-count}
Show that the number of changes in the ordered short-vector support data over an $H$-interval
of length $T$ is $o(qT)$, with an explicit defect small enough to force either a stable
subspace or a contradiction with the rank threshold.
\end{target}

This target uses Baker theory in a role different from direct approximation: not to exclude a
single counterexample, but to limit how rapidly the lattice spectrum can reorganize.

\section{Dilation rigidity and multiscale persistence}
\label{sec:dilation-rigidity}

The normalized orbit sets are nested, while unnormalized points evolve by
\[
 (X,Y)\longmapsto(2X,3Y).
\]
Exact reuse of an auxiliary polynomial across scales is severely rigid.

\subsection{A sparse dilation theorem}

\begin{theorem}[Consecutive-dilation rigidity]
\label{thm:dilation-rigidity}
Let
\[
 P(X,Y)=\sum_{e=(i,j)\in E}c_eX^iY^j
\]
have $q$ nonzero monomials, and let $X,Y>0$.  If
\begin{equation}
 P(2^nX,3^nY)=0
 \qquad(n=n_0,n_0+1,\ldots,n_0+q-1),
 \label{eq:q-dilations}
\end{equation}
then $P=0$.
\end{theorem}

\begin{proof}
Equation~\eqref{eq:q-dilations} is the square linear system
\[
 \sum_{e=(i,j)\in E}(c_eX^iY^j)(2^i3^j)^n=0.
\]
The bases $2^i3^j$ are distinct.  After factoring their $n_0$th powers, the coefficient matrix
is an ordinary Vandermonde matrix and is invertible.  Hence every $c_eX^iY^j$ is zero, so every
$c_e$ is zero.
\end{proof}

\begin{corollary}[No fixed trap for a persistent orbit point]
A nonzero $q$-monomial polynomial cannot vanish at $q$ consecutive unnormalized incarnations
of the same normalized orbit point.
\end{corollary}

This matters because a fixed orbit index $k$ persists: once $k\le K_H$,
\[
 (X_{H+1,k},Y_{H+1,k})=(2X_{H,k},3Y_{H,k}).
\]
A multiscale proof cannot simply discover one polynomial and replay it indefinitely.

\subsection{Invariant subspaces of the dilation operator}

Define
\[
 \mathscr D(P)(X,Y):=P(2X,3Y).
\]
On monomials,
\[
 \mathscr D(X^iY^j)=2^i3^jX^iY^j.
\]
The eigenvalues are simple.

\begin{proposition}[Invariant-space classification]
\label{prop:invariant-spaces}
Every finite-dimensional $\mathscr D$-invariant subspace of $\R[X,Y]$ is spanned by
monomials.
\end{proposition}

\begin{proof}
The operator is diagonal in the monomial basis with pairwise distinct eigenvalues.  Every
invariant subspace of a diagonal operator with simple spectrum is the span of a subset of its
eigenvectors.
\end{proof}

A monomial never vanishes on the positive curve.  Therefore an exactly invariant trap space
cannot solve the problem.

\subsection{The viable replacement: controlled covariance}

Exact invariance is too rigid, but adjacent trap lattices are not unrelated.  The coefficient
weights are multiplied by $2^i3^j$, while the normalized point set changes by at most one
orbit insertion.  A useful theorem would permit a low-rank error:

\begin{target}[Approximate dilation covariance]
\label{target:approx-covariance}
Construct trap subspaces $W_H\subset\R^{E_H}$ and transport maps
$T_H:W_H\to W_{H+1}$ such that
\begin{enumerate}
\item $\dim W_H=(1-o(1))|E_H|$;
\item the uniformly short property is preserved by $T_H$;
\item $\dim(W_{H+1}/T_HW_H)=o(H)$;
\item the accumulated transport defect is incompatible with
Theorem~\ref{thm:dilation-rigidity} unless no nonintegral orbit is present.
\end{enumerate}
\end{target}

The fourth clause is essential: without it, a new basis can be chosen independently at every
scale and nesting supplies no leverage.

\section{Stable factors and unit-equation degeneration}
\label{sec:stable-factors}

If two short polynomials share a factor, the general Brisebarre--Hanrot algorithm may return
``fail.''  On the power curve, common factors are not an obstacle to point counting, but they
may contain useful multiscale information.

\subsection{Finite recurrence is impossible}

\begin{proposition}[Finite projective recurrence]
\label{prop:finite-recurrence}
Let $V\subset\R[X,Y]$ be finite dimensional.  Suppose that for infinitely many $H$ there is a
nonzero $P_H\in V$ vanishing on the entire normalized set $\cC_H$, and suppose the projective
classes $[P_H]\in\mathbb P(V)$ take only finitely many values.  Then no counterexample exists.
\end{proposition}

\begin{proof}
One projective class occurs for infinitely many $H$.  Since the sets $\cC_H$ are nested and
their cardinalities tend to infinity under a counterexample, the corresponding fixed
polynomial vanishes at infinitely many points of the power arc.  Corollary
\ref{cor:sparse-zero-bound} forces it to be zero, a contradiction.
\end{proof}

The proposition remains true when ``finitely many values'' is replaced by any condition that
forces a fixed nonzero limit polynomial and controls convergence strongly enough to pass zeros
to the limit.

\subsection{Where unit-equation theory may enter}

Factor coefficients produced by trap lattices are large integers with structured prime
content.  If a bounded-degree factor recurs after the dilation rescaling
$(X,Y)\mapsto(2X,3Y)$, its coefficients satisfy exponential recurrences.  Relations among
several such factors become one-parameter families of unit equations.  The 2026 results of
Wang--Xiao show that, under transverse boundary-divisor hypotheses, such families are forced
into proper algebraic subsets and yield new gcd estimates for exponential Diophantine
problems~\cite{WangXiao2026}.

The missing bridge is explicit:

\begin{target}[Stable-factor degeneration]
\label{target:stable-factor}
Show that a common factor occurring in a positive-density set of optimized trap spaces must,
after monomial normalization, belong to a fixed algebraic family whose coefficient evolution
is governed by a unit equation.  Prove that every degenerate family is monomial or binomial.
The monomial case has no positive zeros; the binomial case would force a rational relation
between $\log2$ and $\log3$.
\end{target}

This route is secondary to the successive-minimum target, but it is the correct response if
large common factors repeatedly appear in exact computations.

\section{Hermite multiplicity: a tempting route and its row-reset obstruction}
\label{sec:jets}

Multiplicity appears to offer an immediate improvement.  If an integer polynomial and all its
partial derivatives of total order below $s$ are uniformly smaller than one, then all those
integer values vanish at every exact point.  Each point becomes a zero of multiplicity at least
$s$ on the restricted exponential polynomial.

\begin{proposition}[Hermite trap bound]
\label{prop:hermite-trap}
Let $P\in\Z[X,Y]$ have $q$ nonzero monomials.  Suppose that for every
$\alpha=(\alpha_1,\alpha_2)\in\Nzero^2$ with $|\alpha|<s$,
\begin{equation}
 \sup_{t\in J}
 \left|\partial_X^{\alpha_1}\partial_Y^{\alpha_2}
 P(2^{H+t},3^{H+t})\right|<1.
 \label{eq:jet-short}
\end{equation}
Then the segment contains at most $(q-1)/s$ exact integer points.
\end{proposition}

\begin{proof}
Every derivative value in~\eqref{eq:jet-short} is an integer at an integer point, hence zero.
Thus $P$ has algebraic multiplicity at least $s$ there.  Composing with the analytic curve
preserves at least that order of vanishing.  Corollary~\ref{cor:sparse-zero-bound} bounds the
total zero multiplicity by $q-1$.
\end{proof}

If $s$ could grow like the square root of $q$ at negligible cost, this would reduce a quadratic
coefficient budget to a linear point bound and might prove the conjecture.  The cost is not
negligible.

\subsection{Why derivatives are not free}

For one derivative family, high-order Chebyshev modes carry the decay $\rho^{-n}$.  Adding all
jets of order below $s$ creates
\[
 J_s=\binom{s+1}{2}
\]
new families, each with its own low Chebyshev modes.  A Cauchy--Binet minor can replace a very
expensive mode $n\asymp q$ by mode $0$ from a derivative family.  The derivative contributes a
scale penalty roughly
\[
 2^{-H\alpha_1}3^{-H\alpha_2}
\]
but saves a Chebyshev penalty $\rho^{-n}$.  In the relevant regime,
$q\log\rho\asymp H^2/\log H$ is much larger than $H$, so these row resets dominate the naive
entrywise smallness of derivatives.

There is also an exact algebraic calibration.

\begin{lemma}[Jet matrix nondegeneracy]
\label{lem:jet-matrix}
Let
\[
 \mathcal T_s:=\{(i,j)\in\Nzero^2:i+j<s\}.
\]
At any point $(X_0,Y_0)$ with $X_0Y_0\ne0$, the square matrix
\[
 \left(\partial_X^{\alpha_1}\partial_Y^{\alpha_2}X^iY^j
       \big|_{(X_0,Y_0)}\right)_{
       \alpha\in\mathcal T_s,\ (i,j)\in\mathcal T_s}
\]
is invertible.
\end{lemma}

\begin{proof}
If a polynomial of total degree below $s$ has all jets of total order below $s$ equal to zero
at $(X_0,Y_0)$, its Taylor polynomial there is identically zero; hence the polynomial is zero.
Thus the jet evaluation map is injective between vector spaces of the same dimension.
\end{proof}

Lemma~\ref{lem:jet-matrix} says that $J_s$ zero-mode jet rows can consume $J_s$ coefficient
directions before any Chebyshev decay is used.  Multiplicity has a quadratic interpolation
tax.

\begin{target}[Subcritical Hermite regime]
\label{target:subcritical-hermite}
Determine the largest $s=s(q,H)$ for which the combined jet--Chebyshev lattice still has a
short nonzero vector.  The interesting range is
\[
 \frac{H}{(\log H)^2}\ll s\ll\frac{H}{\log H},
\]
where multiplicity would reduce $q/s$ to $o(H)$ but the jet count
$J_s$ is still $o(q)$.  A positive result in this range would prove the conjecture; a matching
negative result would close the most direct multiplicity escape.
\end{target}

The target is concrete: it is a min-plus assignment problem between derivative order and
Chebyshev mode, followed by one covolume estimate.  No transcendence theorem is hidden in its
statement.

\section{Exact Bessel factorization of the analytic lattice}
\label{sec:bessel}

The generic Bernstein-ellipse estimate is sufficient for asymptotics, but the present curve
has an exact Chebyshev expansion.  With $u=2t-1$,
\[
 e^{\lambda(H+t)}=e^{\lambda(H+1/2)}e^{(\lambda/2)u}.
\]
The classical generating formula for modified Bessel functions is
\begin{equation}
 e^{zu}=I_0(z)+2\sum_{n\ge1}I_n(z)T_n(u).
 \label{eq:bessel-cheb}
\end{equation}
Therefore
\begin{equation}
 a_{n,e}=\epsilon_n e^{\lambda_e(H+1/2)}I_n(\lambda_e/2),
 \qquad
 \epsilon_0=1,\quad\epsilon_n=2\ (n\ge1).
 \label{eq:exact-cheb-coeff}
\end{equation}

\begin{proposition}[Diagonal-flow factorization]
\label{prop:bessel-factorization}
For the first $N$ Chebyshev rows and a support $E$, the coefficient matrix factors as
\begin{equation}
 A_N(H,E)=D_{\rm row}\,K_N(E)\,D_H,
 \label{eq:bessel-factorization}
\end{equation}
where
\[
 (D_{\rm row})_{n,n}=\epsilon_n,
 \qquad
 K_N(E)_{n,e}=I_n(\lambda_e/2),
 \qquad
 (D_H)_{e,e}=e^{\lambda_e(H+1/2)}.
\]
All dependence on $H$ is in the final diagonal matrix.
\end{proposition}

This makes the parametric-geometry structure exact rather than heuristic.  It also gives a
clean finite-verification route: interval enclosures for $\log2$, $\log3$, $e^x$, and
$I_n(x)$ produce certified enclosures for every matrix entry, while the Bessel series gives a
positive tail bound.

\subsection{Compound matrices as the spectral bottleneck}

The $k$th exterior power satisfies
\[
 \bigwedge^k A_N(H,E)=
 (\bigwedge^kD_{\rm row})(\bigwedge^kK_N(E))(\bigwedge^kD_H).
\]
The last factor is explicit: its diagonal entries are
$e^{H\sum_{e\in F}\lambda_e}$ over $k$-subsets $F\subset E$.  Thus every nontrivial
condition-number question is concentrated in the fixed Bessel matrix $K_N(E)$.

\begin{target}[Bessel compound bounds]
\label{target:bessel-compounds}
For weighted-simplex frequency sets, obtain uniform lower and upper bounds for a
$1-o(1)$ fraction of the singular values of $K_q(E)$, or equivalently for the bulk of its
compound minors.  Bounds strong enough to prevent concentration of covolume in $o(q)$ singular
directions would convert the determinant estimate into the high-corank conclusion required by
Corollary~\ref{cor:rank-frontier}.
\end{target}

The entries of $K$ are positive and the underlying exponential kernel is strictly totally
positive.  This makes total-positivity techniques, oscillatory matrices, and variation-
diminishing factorizations natural tools.  The report does not assume that the finite Bessel
matrix is uniformly well conditioned; establishing or refuting precisely that assertion is
the point of Target~\ref{target:bessel-compounds}.

\subsection{An exact tail formula}

For $z\ge0$, every $I_n(z)$ is positive and
\[
 I_n(z)=\sum_{k\ge0}\frac{(z/2)^{n+2k}}{k!(n+k)!}.
\]
Hence a rational interval upper bound for $z$ gives a monotone, proof-producing upper bound on
\[
 2e^{\lambda(H+1/2)}\sum_{n\ge N}I_n(\lambda/2).
\]
One may bound this series either termwise or by the Bernstein estimate
\eqref{eq:tail-weight}.  In a finite certificate, the smaller of the two bounds can be stored.
No numerical quadrature is required.

\section{General multiplicatively independent bases}
\label{sec:general-bases}

The method is not special to $2$ and $3$.  Let $u,v>1$ be multiplicatively independent
integers and put
\[
 a_u=\log u,
 \qquad b_v=\log v.
\]
Define
\[
 I_{u,v}(H):=\#\{x\in[H,H+1):u^x,v^x\in\Z\}.
\]
Every monomial restricts to the frequency $i\log u+j\log v$, and the exact block scales
$u^H,v^H$ are integers.

\begin{theorem}[Two-base block theorem]
\label{thm:general-bases}
For multiplicatively independent integers $u,v>1$,
\begin{equation}
 \limsup_{H\to\infty}
 I_{u,v}(H)\frac{(\log H)^2}{H^2}
 \le\frac{32\log u\log v}{9}.
 \label{eq:general-base-bound}
\end{equation}
\end{theorem}

\begin{proof}
Repeat Sections~\ref{sec:total-positivity}--\ref{sec:block-bound} with frequencies
$i\log u+j\log v$ and the weighted simplex
$i\log u+j\log v\le L$.
\end{proof}

The constant is smallest when the logarithms are unbalanced, but the conjectural structure of
counterexamples also changes.  The theorem is best viewed as a general sparse-curve counting
principle rather than evidence that one pair of bases is intrinsically easier.

\section{Proof-producing finite certificates}
\label{sec:certificate-design}

The asymptotic theorem is useful conceptually, but a finite verifier can extend zero-free
regions and provide data about the full successive-minimum spectrum.  The essential principle
is separation of \emph{discovery} from \emph{verification}: LLL, floating-point Bessel values,
and support heuristics may discover a candidate; the final certificate is checked by exact
integer arithmetic plus outward-rounded intervals.

\subsection{Certificate contents}

For each interval $J=[\alpha,\beta]\subset[0,1]$, store:

\begin{enumerate}
\item the integer scale $H$ and rational endpoints $\alpha,\beta$;
\item the ordered monomial support $E$;
\item integer coefficient vectors $c^{(1)},\ldots,c^{(r)}$;
\item an exact nonzero $r\times r$ minor certifying their rank;
\item rational intervals enclosing the first $N$ Chebyshev coefficients of every restricted
monomial;
\item rational upper bounds for all tails;
\item a rational number $U_a<1$ satisfying
\[
 \sum_{n<N}\left|\sum_ec^{(a)}_e a_{n,e}\right|
 +\sum_e|c^{(a)}_e|\operatorname{Tail}_e\le U_a.
\]
\end{enumerate}

The verifier never needs to reproduce the LLL reduction.  It checks the integer minor, the
coefficient enclosures, and the displayed rational inequality.  Theorem
\ref{thm:rank-codim} then returns the point bound $|E|-r$.

\subsection{Directed rounding}

The only transcendental constants are $\log2$, $\log3$, exponentials, and Bessel values at
positive arguments.  All admit rapidly convergent monotone series or standard interval-library
implementations.  A robust verifier can use either:

\begin{itemize}
\item MPFR/Arb directed rounding to produce rational dyadic endpoints; or
\item a small standalone rational checker for truncated Taylor/Bessel series with a proved
remainder.
\end{itemize}

The second option is slower but minimizes the trusted code base and is well suited to Lean
replay.

\subsection{Discovery algorithm}

A practical search can use the following pipeline.

\begin{lstlisting}[style=pseudocode]
INPUT: scale H, interval J, target defect d_target
1. Choose a weighted-frequency support E and Chebyshev cutoff N >= |E|.
2. Compute high-precision enclosures for the Bessel coefficient matrix.
3. Build an integer-scaled approximation lattice, including tail rows.
4. Run LLL/BKZ to obtain many candidate coefficient vectors.
5. Sort candidates by rigorously verified curve sup bound.
6. Retain vectors with bound < 1 and compute their exact Z-rank.
7. If |E| - rank <= d_target, emit a certificate for J.
8. Otherwise split J or modify E, N, and the lattice scaling.
OUTPUT: exact certificate or a diagnostic successive-minimum profile.
\end{lstlisting}

For finite exclusion of a nonintegral point in a unit block, one can account separately for
the known integral point $(2^H,3^H)$.  A cover defect of at most one proves that there is no
additional point.  A defect of zero on the open interval $(0,1)$ is even cleaner.

\subsection{Why this can outperform point scanning}

A direct scan examines candidate integers one at a time.  A trap certificate excludes an
entire continuum interval at once.  The cost is shifted from the number of possible $X$ values
to the dimension of a structured lattice, which grows polynomially in $H$.  The asymptotic
theorem does not yet make this polynomial small enough for a proof, but for finite $H$ the
constant and the observed rank can be far better than the worst-case estimates.

\begin{experiment}[Certificate reconnaissance]
\label{exp:recon}
\statuscompute\quad
For each $14\le H\le40$, compute the full verified rank curve
\[
 r\longmapsto\mu_r(H)
\]
for several weighted-simplex and thin-shell supports.  Record:
\begin{enumerate}
\item the certified defect $q-r$ at threshold;
\item the distribution of coefficient supports and signs;
\item common factors among short polynomials;
\item the change under $H\mapsto H+1$;
\item the singular spectrum of the fixed Bessel matrix.
\end{enumerate}
The aim is not merely to enlarge a finite bound, but to decide empirically which of
Targets~\ref{target:trap-spectrum}, \ref{target:bessel-compounds},
\ref{target:approx-covariance}, or \ref{target:stable-factor} has the right shape.
\end{experiment}

\section{Lean formalization plan}
\label{sec:lean-plan}

The core deterministic theorem is unusually formalization-friendly.  It separates into small
modules with few analytic dependencies.

\subsection{Phase I: exact algebra}

\begin{enumerate}
\item \texttt{ExpFrequency}: prove injectivity of
$(i,j)\mapsto i\log2+j\log3$ from unique factorization.
\item \texttt{ExpVandermonde}: formalize rank of
$(e^{\lambda_jt_i})$ for ordered distinct nodes and frequencies.  A Rolle-theorem proof avoids
heavy total-positivity libraries.
\item \texttt{CurveEvaluationRank}: package Corollary~\ref{cor:evaluation-rank}.
\item \texttt{RankCodimensionTrap}: prove Theorem~\ref{thm:rank-codim} from integer-valuedness
and matrix rank.
\item \texttt{DilationRigidity}: formalize Theorem~\ref{thm:dilation-rigidity} as an ordinary
Vandermonde argument over $\R$.
\end{enumerate}

\subsection{Phase II: certificate checker}

The checker need not formalize LLL or Minkowski.  Its input is a finite list of integer vectors
and rational interval bounds.  Required modules are:

\begin{enumerate}
\item rational interval arithmetic;
\item exact determinant/rank checking over $\Z$;
\item a verified bound for Chebyshev polynomials on $[-1,1]$;
\item a verified exponential/Bessel tail inequality;
\item interval cover bookkeeping.
\end{enumerate}

This phase can certify concrete finite exclusions independently of the asymptotic existence
proof.

\subsection{Phase III: asymptotic theorem}

Formalizing Theorem~\ref{thm:block-bound} requires:

\begin{enumerate}
\item Bernstein-ellipse coefficient bounds;
\item Cauchy--Binet and Hadamard estimates;
\item Minkowski's second theorem for Euclidean lattices;
\item the weighted-simplex lattice-point asymptotics;
\item routine logarithmic asymptotics for the choice~\eqref{eq:Lchoice}.
\end{enumerate}

This is substantial but modular.  The recommended order is to formalize the rank--codimension
checker first, because it can validate computational certificates long before the general
Minkowski theorem is available in the repository.

\section{Comparison with the inherited methods}
\label{sec:comparison}

\begin{longtable}{@{}>{\raggedright\arraybackslash}p{0.19\textwidth}>{\raggedright\arraybackslash}p{0.26\textwidth}>{\raggedright\arraybackslash}p{0.21\textwidth}>{\raggedright\arraybackslash}p{0.19\textwidth}@{}}
\toprule
Method & Arithmetic object & Best scale exposed & Remaining obstruction \\
\midrule
\endfirsthead
\toprule
Method & Arithmetic object & Best scale exposed & Remaining obstruction \\
\midrule
\endhead
Four exponentials & $2\times2$ exponential matrix & Exact qualitative barrier & Conjectural transcendence theorem \\
Determinant/$p$-adic methods in Report IV & Integer semigroup determinants & Sharp termwise and cancellation ceilings & Missing second source of size \\
O-minimal counting & Rational points on fixed arc & Polylogarithmic upper bounds & Exponent too large; denominator support ignored \\
Brisebarre--Hanrot general algorithm & Two small auxiliary polynomials & Efficient local point isolation & Possible common factor \\
This report: one sparse trap & Exponential Chebyshev sum & $H^2/(\log H)^2$ block bound & Still superlinear \\
This report: rank--codimension & Full short-vector subspace & Exact defect $q-r$ & Need $>98.59\%$ rank \\
This report: Bessel/PGN & Successive-minimum spectrum & Exact diagonal-flow factorization & Bulk singular-value theorem missing \\
Hermite jets & Partial-derivative trap & Potential multiplicity gain & Quadratic jet-row reset tax \\
Stable factors & Dilation-covariant factor families & Could convert failure into rigidity & Unit-equation bridge not yet built \\
\bottomrule
\end{longtable}

The new route is complementary to Report IV's determinant program.  Both manufacture small
integer-valued analytic objects, but they organize arithmetic differently.  The previous
program studies determinants of values at many semigroup points; the present program studies
many coefficient vectors small on an entire curve segment.  Their shared barrier is now
visible as a rank statement: local volume decay is easy, while broad, correlated cancellation
across almost all directions is hard.

\section{Prioritized research program}
\label{sec:research-program}

The calculations above narrow the next work to five testable programs.  They are ordered by
expected information gain, not by apparent elegance.

\subsection{Priority A: build the deterministic finite verifier}

Implement Section~\ref{sec:certificate-design} first.  This has three immediate payoffs:

\begin{enumerate}
\item it can extend the finite zero-free region without enumerating every integer candidate;
\item it measures the actual rank at the strict threshold, which the asymptotic covolume cannot
predict;
\item it reveals whether common factors, stable supports, or Bessel condition numbers are the
real finite-scale phenomenon.
\end{enumerate}

\textbf{Kill criterion.}  If optimized supports through the feasible range consistently have
rank defect comparable to $q$, then near-full-rank closure is implausible without new
arithmetic rows.  If the defect is small or decreases rapidly with $H$, the PGN program becomes
the highest priority.

\subsection{Priority B: analyze the fixed Bessel matrix}

The factorization~\eqref{eq:bessel-factorization} isolates the only non-diagonal analytic
object.  Compute and prove:

\begin{itemize}
\item signs and magnitudes of all contiguous minors;
\item asymptotic singular-value density for weighted-simplex frequencies;
\item variation-diminishing or Neville-elimination factorizations;
\item compound-matrix bounds uniform in $q$.
\end{itemize}

\textbf{Kill criterion.}  An exponentially broad singular-value distribution with a positive
fraction above the trap threshold would show that analytic total positivity alone cannot meet
the high-corank target.

\subsection{Priority C: add arithmetic rows tied to a primitive pair}

The unconditional lattice sees only the curve and therefore treats the integral branch and a
hypothetical second generator symmetrically.  A closing proof must likely add rows that exist
only when integer coordinates have the fixed semigroup form
$(2^nM^k,3^nA^k)$.  Candidate rows include:

\begin{itemize}
\item congruences at structural primes dividing $MA$ but not $6$;
\item common-factor data from adjacent normalized orbit points;
\item exact valuation constraints on coefficients evaluated over a whole nested block;
\item transported factors under $(X,Y)\mapsto(2X,3Y)$.
\end{itemize}

The objective is not more covolume in a few directions.  It is to lift the bottom of the
successive-minimum spectrum, so that a determinant bound forces the remaining directions to be
short.

\textbf{Kill criterion.}  If every added arithmetic row is supported on only $o(q)$ coefficient
directions or contributes only $O(H)$ total logarithmic weight, it cannot alter the
$98.59\%$ rank requirement.

\subsection{Priority D: solve the jet min-plus problem}

Target~\ref{target:subcritical-hermite} is the only route in this report that could change the
point-count exponent directly.  Form the combined derivative--Chebyshev matrix, assign to a row
$(\alpha,n)$ the leading cost
\[
 H(\alpha_1\log2+\alpha_2\log3)+n\log\rho,
\]
and compute the minimum total cost of a full-rank row selection.  The exact falling-factorial
minors decide which assignments are admissible.

\textbf{Positive outcome.}  A multiplicity
$s\gg H/(\log H)^2$ with a support of size $O(H^2/(\log H)^2)$ proves the conjecture.

\textbf{Negative outcome.}  A matching assignment lower bound showing
$s=O(H/(\log H)^2)$ closes the most direct Hermite escape and supplies a reusable no-go theorem.

\subsection{Priority E: stable factors and unit equations}

Run factorization on every certified short polynomial and track factors after dilation
normalization.  If a bounded-complexity factor family recurs, attempt the Wang--Xiao
integral-point degeneration route.  If factors do not recur, the data supports a pure PGN
model instead.

\section{Conclusion}

The attached report showed that the conjecture lies exactly on the four-exponentials frontier
and closed many plausible determinant, $p$-adic, cyclotomic, and dynamical shortcuts.  The new
2026 Coppersmith--Chebyshev input does not magically cross that frontier, but it changes the
shape of the search in four concrete ways.

First, the problem is now an exact rational-grid problem on one fixed analytic arc.  Second,
strict total positivity removes the generic coprimality heuristic and turns every independent
short vector into one unit of point-count codimension.  Third, optimizing the frequency support
and the Chebyshev ellipse gives a new unconditional logarithmic-square saving:
\[
 I(H)\le(2.7075555926\ldots+o(1))\frac{H^2}{(\log H)^2}.
\]
Fourth, optimizing over macroscopic windows converts the residual gap into a precise spectral
target: at the present kernel frontier, more than $98.59\%$ of the coefficient space must be
certified short.

That target is difficult, but it is neither vague nor equivalent to simply restating the
conjecture.  It can be attacked computationally, through Bessel compound matrices, through
parametric geometry of numbers, through structural-prime rows, or through controlled
multiscale covariance.  Each attack has an exact success criterion and an exact failure mode.
The next decisive experiment is therefore clear: build the proof-producing verifier, measure
the full successive-minimum spectrum, and let the observed defect determine which theorem to
pursue.

\appendix

\section{A Rolle-theorem proof of the sparse zero bound}
\label{app:rolle}

For formalization it is useful to avoid invoking the general theory of total positivity.
We record the elementary induction.

\begin{theorem}[Exponential sums have few zeros]
Let $\lambda_1<\cdots<\lambda_q$ and
\[
 F(t)=\sum_{j=1}^q c_je^{\lambda_jt}
\]
be nonzero.  Then $F$ has at most $q-1$ real zeros counted with multiplicity.
\end{theorem}

\begin{proof}
Induct on $q$.  The assertion is clear for $q=1$.  Suppose it holds for $q-1$ and divide by
the first exponential:
\[
 G(t)=e^{-\lambda_1t}F(t)
     =c_1+\sum_{j=2}^q c_je^{(\lambda_j-\lambda_1)t}.
\]
If $G$ has $m$ zeros counted with multiplicity, then $G'$ has at least $m-1$ zeros counted with
multiplicity.  This follows by applying Rolle between distinct zeros and reducing each positive
multiplicity by one at the same location.  But
\[
 G'(t)=\sum_{j=2}^q c_j(\lambda_j-\lambda_1)e^{(\lambda_j-\lambda_1)t}
\]
is either zero or a nontrivial sum of at most $q-1$ exponentials.  If it is nonzero, induction
gives $m-1\le q-2$.  If it is zero, then all $c_j$ for $j\ge2$ vanish and $G$ is a nonzero
constant, so $m=0$.  Thus $m\le q-1$.
\end{proof}

The same induction proves nonsingularity of every square evaluation matrix: a nonzero kernel
vector would give a nontrivial sum with at least as many distinct zeros as terms.

\section{Detailed fixed-rank successive-minimum estimate}
\label{app:fixed-rank}

Let $B$ be the matrix~\eqref{eq:Bmatrix}, let $\Delta=\sqrt{\det(B^TB)}$, and let
$\mu_1\le\cdots\le\mu_q$ be the Euclidean successive minima.  The volume of the unit ball is
\[
 \omega_q=\frac{\pi^{q/2}}{\Gamma(q/2+1)},
 \qquad
 -\log\omega_q=\frac q2\log q+O(q).
\]
Minkowski gives
\begin{equation}
 \mu_1\cdots\mu_q\le\frac{2^q\Delta}{\omega_q}.
 \label{eq:minkowski-product}
\end{equation}
Because $(0,0)\in E$, the smallest diagonal tail weight is
\[
 \tau_{0,0}=\frac{2\rho^{-q}}{1-\rho^{-1}}\ge2\rho^{-q}.
\]
For every nonzero $c\in\Z^E$,
$\|Bc\|_2\ge2\rho^{-q}$.  Hence, for fixed $r$,
\[
 (2\rho^{-q})^{r-1}\mu_r^{q-r+1}
 \le\mu_1\cdots\mu_q.
\]
Together with~\eqref{eq:minkowski-product},
\begin{align*}
 \log\mu_r
 &\le\frac{1}{q-r+1}
 \left(\log\Delta+(r-1)q\log\rho+O(q\log q)\right)\\
 &\le\frac{H+R}{q-r+1}\sum_e\lambda_e
 -\frac{q(q-1)}{2(q-r+1)}\log\rho
 +\frac{(r-1)q}{q-r+1}\log\rho+O_r(\log q).
\end{align*}
For fixed $r$, replacing $q-r+1$ by $q$ in the leading terms changes the expression by
$O_r((H\bar\lambda)/q+\log\rho)$, which is lower order in the applications.  This yields
\eqref{eq:fixed-r-minimum}.

\section{Numerical constants and optimization data}
\label{app:constants}

All logarithms in this report are natural.  The constants are:
\begin{align*}
 \log2&=0.6931471805599453094\ldots,\\
 \log3&=1.0986122886681096914\ldots,\\
 \frac{32\log2\log3}{9}&=2.707555592600207\ldots.
\end{align*}
The root of $t\tanh t=2$ is
\[
 t_0=2.065338138974704728\ldots,
 \qquad
 \rho_0=e^{t_0}=7.887964675046528\ldots.
\]
Then
\begin{align*}
 R(\rho_0)&=2.503685361986428\ldots,\\
 r_0&=\frac1{R(\rho_0)-1}=0.665032892479158\ldots,\\
 \kappa_*&=5.087266580727214\ldots.
\end{align*}
The rank thresholds are summarized below.

\begin{center}
\small\begin{tabular}{@{}>{\raggedright\arraybackslash}p{0.23\textwidth}>{\raggedright\arraybackslash}p{0.17\textwidth}>{\raggedright\arraybackslash}p{0.23\textwidth}>{\raggedright\arraybackslash}p{0.20\textwidth}@{}}
\toprule
Input on least generator & $\kappa_*\beta$ & Required rank fraction & Maximum defect fraction \\
\midrule
Kernel level $\beta\ge14$ & $71.221732\ldots$ & $0.98595934\ldots$ & $0.01404066\ldots$ \\
Software level $\beta>27$ & $137.356198\ldots$ & $0.99271966\ldots$ & $0.00728034\ldots$ \\
\bottomrule
\end{tabular}
\end{center}

For comparison, using the unweighted total-degree triangle would give the larger window
constant
\[
 \frac{16(\log6)^2}{9}\frac{\cosh t_0}{t_0^2}
 =5.361841943267\ldots.
\]
The weighted simplex reduces it by the exact factor
$4\log2\log3/(\log6)^2=0.948\ldots$.

\section{A machine-checkable certificate schema}
\label{app:schema}

A compact textual schema can be organized as follows.

\begin{lstlisting}[style=pseudocode]
certificate_version: 1
scale_H: integer
intervals:
  - left: rational
    right: rational
    support: [[i_1,j_1],...,[i_q,j_q]]
    chebyshev_cutoff: N
    polynomials:
      - coefficients: [c_1,...,c_q]
        rank_minor_row: optional
        verified_upper_bound: rational < 1
    rank_minor:
      row_indices: [...]
      column_indices: [...]
      determinant: nonzero integer
    coefficient_intervals:
      # interval enclosure for each n,e
    tail_bounds:
      # nonnegative rational for each e
claimed_point_bound: sum(q-r)
\end{lstlisting}

A checker performs no search.  It verifies:

\begin{enumerate}
\item every support exponent is nonnegative;
\item coefficient vectors and the stored minor agree;
\item the minor determinant is nonzero;
\item interval arithmetic proves each displayed upper bound;
\item intervals cover the requested parameter range;
\item the claimed point bound equals the sum of local defects.
\end{enumerate}

The mathematical kernel is Theorem~\ref{thm:rank-codim}; all other operations are finite.

\section{Dependency graph of the new results}
\label{app:dependency}

\begin{center}
\begin{tabularx}{0.96\textwidth}{@{}p{0.27\textwidth}X@{}}
\toprule
Result & Dependencies \\
\midrule
Sparse zero bound & Multiplicative independence of $2,3$; Rolle's theorem \\
Rank--codimension trap & Sparse evaluation rank; integer-valuedness \\
Norm-to-sup transfer & Chebyshev expansion; Cauchy--Schwarz \\
Covolume bound & Bernstein coefficient bound; Cauchy--Binet; Hadamard \\
Fixed-rank trap & Covolume bound; Minkowski's second theorem \\
Weighted block theorem & Fixed-rank trap; weighted simplex count and first moment \\
Window constant & Weighted threshold; elementary two-variable optimization \\
Dilation rigidity & Ordinary Vandermonde determinant \\
Hermite trap & Integer derivatives; sparse zero multiplicity \\
Finite certificate checker & Rank--codimension theorem; rational interval arithmetic \\
\bottomrule
\end{tabularx}
\end{center}

\begin{thebibliography}{99}

\bibitem{AE1944}
L.~Alaoglu and P.~Erd\H{o}s,
\newblock On highly composite and similar numbers,
\newblock \emph{Transactions of the American Mathematical Society} \textbf{56} (1944),
448--469.

\bibitem{Baker1975}
A.~Baker,
\newblock \emph{Transcendental Number Theory},
\newblock Cambridge University Press, 1975.

\bibitem{BrisebarreHanrot2026}
N.~Brisebarre and G.~Hanrot,
\newblock Integer points close to a transcendental curve: an algorithmic approach,
\newblock arXiv:2606.04858, 2026.

\bibitem{Cassels1959}
J.~W.~S.~Cassels,
\newblock \emph{An Introduction to the Geometry of Numbers},
\newblock Springer, 1959.

\bibitem{Karlin1968}
S.~Karlin,
\newblock \emph{Total Positivity}, Vol.~I,
\newblock Stanford University Press, 1968.

\bibitem{Pinkus2010}
A.~Pinkus,
\newblock \emph{Totally Positive Matrices},
\newblock Cambridge University Press, 2010.

\bibitem{ReportIV}
V.~Reshetnikov and the ProveIt project,
\newblock \emph{The Alaoglu--Erd\H{o}s Integer-Exponent Conjecture: Unified Report},
\newblock attached Report IV, 2026.

\bibitem{Roy2015}
D.~Roy,
\newblock On Schmidt and Summerer parametric geometry of numbers,
\newblock \emph{Annals of Mathematics} \textbf{182} (2015), 739--786.

\bibitem{SchmidtSummerer2009}
W.~M.~Schmidt and L.~Summerer,
\newblock Parametric geometry of numbers and applications,
\newblock \emph{Acta Arithmetica} \textbf{140} (2009), 67--91.

\bibitem{SchmidtSummerer2013}
W.~M.~Schmidt and L.~Summerer,
\newblock Diophantine approximation and parametric geometry of numbers,
\newblock \emph{Monatshefte f\"ur Mathematik} \textbf{169} (2013), 51--104.

\bibitem{TaoJacobian2026}
T.~Tao,
\newblock A digestion of the Jacobian conjecture counterexample,
\newblock blog article, 21 July 2026,
\newblock \url{https://terrytao.wordpress.com/2026/07/21/a-digestion-of-the-jacobian-conjecture-counterexample/}.

\bibitem{EpochHadamard2026}
Epoch AI,
\newblock Hadamard matrix of order 668,
\newblock research announcement, 2026,
\newblock \url{https://epoch.ai/frontiermath/open-problems/hadamard}.

\bibitem{Waldschmidt2000}
M.~Waldschmidt,
\newblock \emph{Diophantine Approximation on Linear Algebraic Groups},
\newblock Springer, 2000.

\bibitem{WangXiao2026}
J.~T.-Y.~Wang and Z.~Xiao,
\newblock Families of unit equations and exponential Diophantine problems via integral points,
\newblock arXiv:2604.26497, 2026.

\end{thebibliography}

\end{document}
